\chapter{*智能体大语言模型的认知视角}\label{chap:cognition}


\begin{learningbox}
本章将介绍，对于大语言模型（large language model，LLM）与智能体大语言模型（agentic LLM）的\emph{认知地位}，目前哪些判断有据可依，哪些仍无法作出。具体而言，学完本章后，你将能够：
\begin{itemize}
    \item 区分\emph{行为}、\emph{内部表征}与\emph{机制}三个层面。
    \item 概述大语言模型是否\emph{理解}语言内容这一问题上的争论。
    \item 评估支持和反对大语言模型具有\emph{世界模型}与空间表征的证据，并解释这为何关系到智能体的规划与行动。
    \item 解释对于\emph{心智理论}测试表现背后的表征，我们已经知道什么，以及这些认识与知行鸿沟有何关系。
    \item 描述有关大语言模型\emph{元认知}与\emph{内省}的争论，并解释为什么自我报告只能提供较弱的证据。
    \item 概述科学的\emph{意识}理论如何应用于大语言模型，并思考过度或不足地赋予模型心理状态各自带来的风险。
    \item 在目前证据仍然稀少的情况下，从理论上讨论智能体扩展（记忆、工具、多模态与行动）如何影响上述各个问题。
\end{itemize}
\end{learningbox}

\section{从行为到认知}
第五部介绍了许多大语言模型与智能体大语言模型的例子，它们的表现可以与人类相当，或者以类似人类的方式行动。例如，模型能够通过错误信念测试，复现人类总体上的道德判断，或者给出富有共情的回应。然而，我们也看到，这些行为往往十分脆弱，微小扰动就可能使模型偏离正确方向；模型在文本中\emph{说}的内容，有时也与其作为智能体所\emph{做}的事情严重脱节，这就是知行鸿沟（knowing--doing gap）。我们曾在多处暂时搁置由此自然产生的问题：通过错误信念测试，是否意味着模型具有心智理论（theory of mind，ToM；第~\ref{sec:theoryofmind}~节）？以共情方式回应的模型，是否有所感受（第~\ref{sec:empathy}~节）？智能体大语言模型对自身或另一个智能体的行为进行反思，究竟意味着什么（第~\ref{sec:roletask}~节）？这些表现又如何与人类或其他动物背后的能力相比较？

本章将讨论这些问题，而不在有时相当激烈的争论中选边站。我们关注的重点不再是模型说了什么或做了什么，而是其背后的表征或过程，是否当得起那些用来描述人类认知的名称：理解、思考、世界模型、心智理论、意识，等等。因此，本章的根本问题可以追溯到人工智能成为一个研究领域之初：Alan Turing 曾提出著名的行为测试，用以回答机器能否思考的问题 \citep{turing1950computing}；John Searle 则认为，无论行为取得多大成功，都不足以终结这一争论 \citep{searle1980minds}。如今的新变化，首先在于当前大语言模型已经能够通过图灵测试 \citep{jones2025turing, gibney2025turing}；正如第五部反复展示的，它们还能通过种类繁多的其他行为测试，而这些测试在人类身上与复杂的认知功能相关。这使得争论不再像 Turing 和 Searle 所处的时代那样，只停留在假设层面。其次，可解释性技术（第~\ref{chap:mi}~章）使我们能够在行为数据之外，补充有关\emph{内部表征}（编码了哪些信息）和\emph{机制}（信息如何计算和使用）的证据，为认知研究打开新的空间，增进我们对机器与人类的理解。

在展开讨论之前，需要先澄清几个概念：支持认知主张的证据处于什么层面，能力与表现有何区别，以及每项认知技能应以什么作为分析单位。

\subsection{三个证据层面}\label{sec:levels}
对于本章的问题，我们区分三个可以收集证据、支持机器认知主张的层面：\emph{行为层面}（将模型视为黑箱，测试它会做什么）、\emph{表征层面}（模型的激活编码了哪些信息），以及\emph{机制层面}（这些信息是否在因果意义上得到使用）。\footnote{第~\ref{chap:mi}~章定义了四种可解释性范式：行为范式、归因范式、基于概念的范式和机制范式。这里将前两种范式归入行为层面，第三、第四种范式则分别对应表征层面与机制层面。还需注意，Marr 提出的计算、算法与实现三个分析层次 \citep{marr1982vision}，与此处的划分相似，但并不相同。}各个层面的证据都有优缺点，因此需要整合不同层面的认识。我们以第~\ref{sec:theoryofmind}~节的错误信念测试为例：Maxi 不在场时，他的巧克力被从蓝色柜子移到了绿色柜子。

\paragraph{行为层面} 在行为层面，我们实施测试并对答案评分，如第五部所做的那样。其优点包括适用于任何系统，无论是基础模型还是微调模型、开放还是封闭的大语言模型、具有还是不具有智能体能力；测试结果原则上也可以直接与人类得分比较。主要缺点是，同一个测试答案既可能来自真正的信念追踪，也可能来自记忆（测试题目及答案已经包含在训练数据中），或表面启发式与虚假相关，\citep[相关讨论见][]{hu2025reevaluatingtheorymindevaluation,kouwenhoven2026traces}。对测试形式作出扰动，例如将柜门描述为透明的，可以缩小可能解释的范围，但证据依然是间接的。

\paragraph{表征层面} 在表征层面，我们询问模型内部编码了哪些信息。主要方法是{\em 线性探测（linear probing）}：训练一个小型分类器，根据某一层的激活来预测某种属性（参见第~\ref{chap:mi}~章及 \citet{alain2016understanding, belinkov2022probing}）。这样就能检验某些概念是否存在。对于错误信念测试，我们可以询问：能否将 Maxi {\em 认为}巧克力所在的位置，与巧克力实际所在的位置分别解码出来？如果可以，模型就编码了测试所依赖的关键区别。这里有两点需要谨慎对待，后文会进一步讨论：(1) 足够强大的探测器可能自己学会了任务，而非从模型中将其解码出来 \citep{hewitt2019designing}；(2) 能从模型中解码出的信息，未必被模型使用。

\paragraph{机制层面} 在机制层面，我们通过观察不同干预对模型行为的影响，来回应后一项顾虑。干预的例子包括消融、将一次运行的激活修补到另一次运行中，或沿探测器找到的方向引导激活（第~\ref{chap:mi}~章）。例如，若将所编码的 Maxi 的信念从蓝色改为绿色，就使模型回答他会到绿色柜子中寻找，那么这一表征很可能在模型行为中发挥了{\em 因果}作用。支持认知主张的最强证据来自这第三个层面。然而，正如我们将看到的，当前文献中这样的证据仍然稀少。一个原因是，这类研究要求能够访问模型权重，因而大多局限于小模型与范围狭窄的任务。\footnote{前沿模型开发者自行开展这类研究时是一个例外，\citep[例如][]{templeton2024scaling, lindsey2025biology}。这类研究虽然富有启发，但也带来了能否独立复现的问题。}

因此，本章将更细致地说明如何实现{\em 构念效度（construct validity）}。当模型的行为在扰动下仍然稳健时，相关变量在内部有所表征；若这些表征进而在因果意义上得到使用，我们就有了有力的依据，支持相应能力确实存在，以及对该能力的判断具有有效性。表~\ref{tab:levels}~概括了支持本章所讨论的各项（假定）能力的证据。其中最突出的特点是，各处均以行为层面的证据为主。对于未经额外增强的大语言模型，世界模型的表征与机制证据最强，但也仅限于少数领域，例如游戏棋盘表征。对于智能体大语言模型，目前则大体上完全缺乏表征与机制层面的证据。

\begin{table}
\footnotesize
\begin{subtable}{\textwidth}
\begin{tabularx}{\linewidth}{@{}p{2.1cm}XXX@{}}
\toprule
 & \textbf{行为层面} & \textbf{表征层面} & \textbf{机制层面} \\
\midrule
语言理解（第~\ref{sec:understanding}~节） &
\emph{大量。}基准测试；经过扰动的任务变体与反事实任务变体 \citep{wu2024reasoning} &
\emph{日益增多。}从文本中恢复颜色与方向的概念空间 \citep{abdou2021language, patel2022mapping} &
\emph{稀少。}处理特定狭窄语言任务的回路 \citep{wang2022interpretability} \\
\addlinespace
世界模型（第~\ref{sec:worldmodels}~节） &
\emph{日益增多。}导航、规划与游戏任务 \citep{momennejad2023cogeval, paglieri2024balrog} &
\emph{日益增多。}可解码棋盘状态、空间与时间 \citep{li2023emergent, gurnee2024language} &
\emph{日益增多。}编辑棋盘表征会改变落子 \citep{li2023emergent, nanda2023emergent} \\
\addlinespace
心智理论（第~\ref{sec:tomcognition}~节） &
\emph{大量。}见第~\ref{sec:theoryofmind}~节 &
\emph{稀少。}可解码自身与他人的信念 \citep{zhu2024language, bortoletto2024benchmarking} &
\emph{稀少。}引导信念表征会改变答案 \citep{zhu2024language} \\
\addlinespace
自我知识（第~\ref{sec:metacognition}~节） &
\emph{日益增多。}校准、自我预测 \citep{kadavath2022language, binder2024looking} &
\emph{稀少。}真实性方向 \citep{azaria2023internal, burns2023discovering} &
\emph{稀少。}概念注入；忠实性测试 \citep{lindsey2025introspection, turpin2023language} \\
\addlinespace
意识（第~\ref{sec:consciousness}~节） &
\emph{稀少。}自我报告，通常被认为证据较弱 \citep{perez2023towards} &
\emph{几乎没有。}没有公认的测度。依据从意识理论导出的指标属性评估架构 \citep{butlin2023consciousness} &
\emph{稀少。}消融 J 空间的激活，为其在高阶任务中的功能性参与提供依据 \citep{gurnee2026verbalize}\\
\bottomrule
\end{tabularx}
\caption{未经额外增强的大语言模型}\label{tab:levels-bare}
\end{subtable}

\bigskip

\begin{subtable}{\textwidth}
\begin{tabularx}{\linewidth}{@{}p{2.1cm}XXX@{}}
\toprule
 & \textbf{行为层面} & \textbf{表征层面} & \textbf{机制层面} \\
\midrule
语言理解（第~\ref{sec:understanding}~节） &
\emph{没有。} &
--- &
--- \\
\addlinespace
世界模型（第~\ref{sec:worldmodels}~节） &
\emph{稀少。}在规划器内部以大语言模型充当模拟器 \citep{hao2023reasoning,gu2024your} &
--- &
--- \\
\addlinespace
心智理论（第~\ref{sec:tomcognition}~节） &
\emph{稀少。}协调博弈中的支撑架构（第~\ref{chap:role}~章）；未在交互过程中实施干预 &
--- &
--- \\
\addlinespace
自我知识（第~\ref{sec:metacognition}~节） &
\emph{稀少。}未经验证的完成报告 \citep{shapira2026agentschaos}；区分测试与部署情境 \citep{needham2025large} &
--- &
--- \\
\addlinespace
意识（第~\ref{sec:consciousness}~节） &
\emph{几乎没有。}智能体场景中的自我报告与任务偏好 \citep{gray2025agentic} &
\emph{几乎没有。}依据全局工作空间理论评估语言智能体架构 \citep{goldstein2024case} &
--- \\
\bottomrule
\end{tabularx}
\caption{在交互中研究的智能体大语言模型}\label{tab:levels-agentic}
\end{subtable}

\caption{本章各个主题在行为、表征与机制层面的证据来源，分别列出 (1)~未经额外增强的大语言模型和 (2)~智能体大语言模型。表中说明了现有研究的丰富程度，并列出少数代表性研究。对于智能体大语言模型，在撰写本书时可获得的证据总体上仍然稀少。}\label{tab:levels}
\end{table}

\subsection{能力与表现} \label{sec:competence}
成功完成一次测试，并不能确立某项能力的存在；失败也不能证明它不存在。认知科学区分了\emph{能力（competence）}与\emph{表现（performance）}：前者指背后的能力，后者指这种能力在特定场合、特定约束下如何呈现 \citep{chomsky1965aspects}。\citet{firestone2020perf} 认为，比较人类与机器时应当尊重这一区分。例如，一个儿童可能因为不理解问题的措辞而无法通过错误信念测试。同样，大语言模型的表现也可能取决于提示格式、词元化、上下文长度与输出约束等因素。某个基础模型即使无法按照某项错误信念测试要求的格式作答，也可能通过它赋予不同候选续写的概率，展现相应的能力 \citep{kouwenhoven2026traces}。比较还可能存在不对等，例如，人类参与者获得了指导、练习试次与注意力检查，模型却接受零样本测试 \citep{lampinen2024can}。

能力与表现的区分，有助于避免两种相反的错误。一种是\emph{拟人化（anthropomorphism）}：过于轻易地从模型的输出和动作中读出类似人类的倾向。计算机是社会行动者（CASA）范式的例子表明（第~\ref{sec:politeness}~节），这种倾向在研究者中也根深蒂固。另一种则是灵长类动物学家 Frans de Waal 所说的\emph{类人性否认（anthropodenial）}\citep{deWaal1999anthropo}，或 \citet{milliere2026anth} 所称的\emph{人类中心偏差（anthropocentric bias）}：因为某项能力不是以容易辨认为人类式的方式实现的，就忽视或否认它，或者要求模型达到连人类也无法满足的标准。


\subsection{分析单位：模型、系统，还是环境？} \label{sec:unit}
一个反复出现的问题是，应以什么作为分析单位：未经额外增强的大语言模型不具备的能力，在加入智能体控制框架（agentic harness），或置于多智能体环境后，可能会在\emph{系统}层面出现。第~\ref{chap:role}~章讨论过一个例子：在 Hanabi 中加入显式推理与验证步骤后，原本失败的协调变成了合理的对局。

哲学家与认知科学家也曾对人类提出类似观点：认知并非仅是大脑内部发生的事情。由此形成了认知的\emph{4E 观（4-E view）}\citep{noe2009outofourheads, Carney2020foure}：
\begin{enumerate}
    \item 具身性（Embodied）：认知由整个身体及其感觉运动能力塑造，而非只由大脑中的神经活动塑造；
    \item 嵌入性（Embedded）：心理过程置身于特定环境之中，并受到环境的约束；每种环境都有各自的物理、社会与文化属性；
    \item 生成性（Enactive）：知觉与思维在同世界的动态、实时交互中涌现，而非来自被动的内部计算；
    \item 延展性（Extended）：认知过程延伸到环境中，将外部工具、技术以及笔记本、计算器之类的辅助物纳入思考，作为其中主动发挥作用的部分。
\end{enumerate}

按照认知的\emph{延展}观，如果一个笔记本能够被可靠地查阅，它就可以发挥与生物记忆相同的功能，应被视为认知过程与认知系统的一部分 \citep{clark1998extended}。这一观点可以与本章问题所涉及的、未经额外增强的大语言模型和智能体大语言模型之间的区分相比较。在未经额外增强的模型中，工作记忆是上下文窗口，长期知识则存储在权重中，只能通过（进一步）训练来改变。智能体控制框架会加入类似笔记本的外部存储：过去的交互记录、事实与数据，在相关时被检索出来。

这一点贯穿全章，因为不同研究谈论认知时，针对的分析单位并不相同。讨论本章选取的各项认知能力时，我们都会考虑这一因素：理解语言内容（第~\ref{sec:understanding}~节），关于空间环境（第~\ref{sec:worldmodelscog}~节）与社会环境（第~\ref{sec:tomcog}~节）的世界模型，内省与自我知识（第~\ref{sec:metacognition}~节），最后是最具推测性的主题——意识（第~\ref{sec:consciousness}~节）。

\section{理解语言内容}\label{sec:understanding}
第~\ref{chap:behavior}~章表明，大语言模型能够理解许多会话含意。例如，它们能正确推理出：如果一个同事得到的赞扬只有{\em 总是准时}，那么人们对他的评价可能并不特别高，因为这句话的言外之意大致是\emph{说到他，我唯一能想到的优点就是准时}。能够作出这种推理的模型，是否可以说\emph{理解}了自己读写的内容？在这个问题上，学界依然存在深刻分歧。

\subsection{怀疑者的论证}
这一立场的经典论证，见于 Searle 的\emph{中文屋（Chinese Room）}思想实验 \citep{searle1980minds}。一个完全不懂中文的人坐在房间里，手边有一本操纵中文符号的规则手册。问题传进来，他按照规则操作，再传出流畅的答案，然而他什么也不理解。\footnote{在今天看来，我们认为这一思想实验的表述带有倾向性：中文在全球拥有 13 亿母语使用者，而 Searle 的文章将其描述为\emph{无法理解的}语言，隐含了强烈的西方偏见，至少应当对此作出反思。}Searle 因而得出结论：符号操作（句法）并不足以产生意义（语义）。\citet{harnard1990symbol} 将这一问题重新表述为\emph{符号接地问题（symbol-grounding problem）}：如果一个系统中的符号只能用其他符号来定义，那么它的处境就如同一个毫无中文基础的人，试图仅凭一本汉汉词典学习中文。要获得意义，符号终究必须与语言之外的事物相联系，例如对物理世界或社会世界的感知与交互。

\citet{bender-koller2020climbing} 在\emph{章鱼测试（octopus test）}中，将这一推理应用于大语言模型。两个人被困在不同的岛上，通过一根海底电缆通信。一只超级聪明的章鱼窃听电缆，学会了消息中的统计模式，最终冒充其中一个人。在闲聊中，这种欺骗行得通；但当另一座岛上的人请求帮助抵御熊的袭击时，伪装便露馅了。这只章鱼从未接触过这些词语所指向的世界。在作者看来，一个仅用语言形式训练的系统，无论规模变得多大，都无法学会意义。沿着同样的思路，\citet{bender2021dangers} 将大语言模型描述为\emph{随机鹦鹉（stochastic parrot）}：它们依据概率模式将语言形式拼接起来，完全不涉及意义。

\subsection{对怀疑论证的回应}
针对怀疑者的论证，出现了三类回应。第一，意义或许原本就不需要指称。\citet{piantadosi2022meaning} 为\emph{概念角色（conceptual role）}观点辩护：一个概念的意义，在很大程度上存在于它与其他概念的关系之中。人类的许多概念，例如{\em 正义}、{\em 素数}和{\em 夸克}，也并非直接锚定在感知上；赋予这些词语意义的概念间关系，恰恰可以从文本中学到。第二，指称或许无须在每个说话者个体身上实现锚定。我们可以指称 Peano 或仙女座星系，而不必与之有过感知上的接触，因为我们属于一个语言共同体，这些词语在其中有着沿用的历史。\citet{mandelkern2024language} 认为，用这一共同体的产物训练的大语言模型，可能以大致相同的方式继承指称。第三，\citet{pavlick2023symbols} 认为，这归根结底是一个经验问题：大语言模型的表征是否具有锚定理应赋予的结构，应当通过研究这些表征来回答（见表~\ref{tab:levels}~的表征层面一列），而不能只依靠思想实验。

\citet{mitchell2023debate} 总结并综合了这场争论，指出学界存在分歧：2022 年一项针对自然语言处理研究者的调查中，约有一半人同意，仅用文本训练的模型也可以在实质意义上理解语言 \citep{michael2022what}。他们提出，这个问题可能并没有单一答案，大语言模型迫使我们区分不同类型的理解。

\subsection{证据}
已有研究从多个来源提出证据，支持仅凭大规模文本学习就能使概念得到锚定的观点；另一些研究结果则支持相反的看法。

\paragraph{映射世界的结构}
\citet{abdou2021language} 表明，纯文本模型中颜色词表征的几何结构，与人类颜色知觉的几何结构相吻合，尽管这些模型从未见过任何颜色。\citet{patel2022mapping} 向 GPT-3 提供了少量方向词示例，如 {\em left}、{\em north}，它们处于一个以文本呈现的网格世界中。模型能够泛化到未见过的词语，在世界旋转后也能做到；但若将映射关系打乱，就无法泛化。这说明，从文本中学到的概念结构，与其所锚定领域的结构同构。

\paragraph{理解浅薄的迹象}
其他结果表明，这种结构即使存在，也可能与人类理解的方式有所不同。用 {\em A is B} 训练的模型，往往无法推理出 {\em B is A}：它们可以说出某位名人的父亲或母亲，却无法在给定这位父亲或母亲时说出该名人，这种失败被称为\emph{反转诅咒（reversal curse）}\citep{berglund2024reversalcurse}。这一问题涉及存储在权重中的知识。如果两个事实都放在上下文窗口中，模型就能毫无困难地反转它们，因此，这一诅咒反映的是训练时事实的存储方式，而非模型无法领会这种关系是对称的。\citet{mccoy2024embers} 表明，正确率取决于\emph{答案}作为一段文本的概率：当解答是高概率句子时，GPT-4 解读一种简单密码的正确率为 51\%，而当解答是低概率句子时，正确率仅为 13\%。人类也无法免于概率效应；凡是在可预测的句子中漏看过错别字的人，都明白这一点。但是，如果任务得到了稳健意义上的\emph{理解}，我们并不会预期出现如此巨大的差距。推理模型在这类任务上的表现好得多，但其正确率与推理长度仍然取决于输出的概率 \citep{mccoy2024o1}。类似地，在熟悉任务的反事实变体上，例如九进制算术，模型表现会急剧下降 \citep{wu2024reasoning}。最后，\citet{mancoridis2025potemkin} 描述了\emph{波将金式理解（Potemkin understanding）}，其名称来自电影布景中的村庄：它们只有房屋立面，后面并没有真正的房子。这类似于第五部描述的知行鸿沟，如今又出现在语言本身内部：研究者观察到，模型能够正确定义某个概念，例如 ABAB 押韵格式，却随后无法应用它；这种失败方式，不会出现在能够给出该定义的人身上。

总的来看，这些结果符合形式语言能力与功能语言能力的区分 \citep{mahowald2023dissociating}：对语言形式的掌握接近人类水平，同时在实际使用中也确实有所理解，但这种理解并不均衡，而且受训练数据的统计特征制约。

\subsection{智能体视角}
章鱼只能窃听电缆，从中接收信号。多模态大语言模型能够接收图像与音频，智能体大语言模型则能行动和交互。一个智能体编写代码，并在运行时观察代码是否出错；或者删除一个文件，此后便再也无法访问它。在这些情况下，它都从世界中获得了反馈，反馈涉及它对 {\em file} 与 {\em delete} 等词语的使用。这提供了怀疑者所要求的那种感知与行动之间的联系，削弱了章鱼论证的前提。然而，这类反馈大多以文本形式进入模型的上下文窗口；只有在智能体轨迹被用于进一步训练时，反馈才会进入权重。此外，模型不承担身体层面的利害后果：至少到目前为止，模型还不会因为自身的行动而真正“死亡”或“伤害自己”。据我们所知，嵌入智能体系统是否真的会改变上述表征，尚未经过测试。

\subsection{结语}
无论是\emph{随机鹦鹉}还是\emph{完全理解者}，都不符合当前证据。纯文本训练能够产生映射世界某些方面的概念结构，但这种结构的运用方式，仍然对表面统计特征十分敏感。哲学争论部分是关于词语用法的争论；至于这种结构有多稳健，以及行动能否改善它，这些经验问题仍然悬而未决。

\section{世界模型与空间认知}\label{sec:worldmodelscog}\label{sec:worldmodels}
在自然界，要采取恰当的行动，就需要在行动之前预测其后果。对于在世界中行动的人工智能智能体，这一点同样适用；尤其是在动作不可逆时，例如第~\ref{sec:moral}~节所述的那些破坏性动作。
智能行为依赖内部世界模型，并可借其在行动前尝试不同选择，这一想法可以追溯到心理学成为一门科学之初：\citet{craik1943nature} 提出，正是为此，生物体在头脑中携带着外部现实的{\em 缩小模型}。\citet{tolman1948cognitive} 表明，大鼠会形成\emph{认知地图（cognitive map）}：学会走一个迷宫后，它们能够走此前从未走过的捷径。这与强化学习存在相似之处：强化学习中的世界模型，是从经验中学到的环境动力学模型，可用于规划 \citep{ha2018world}。大语言模型是否具有这样的模型，仍有争议。在 \citet{lecun2022path} 看来，自回归大语言模型从原理上就缺乏这类模型，因此他为未来的生成模型提出了不同架构。过去几年里，这一主题已经积累了大量文献。我们将回顾其中几项代表性研究，主要聚焦空间认知，介绍支持和反对大语言模型具有世界模型的证据。

\subsection{支持的证据}
最著名的案例是 Othello-GPT \citep{li2023emergent}，如图~\ref{fig:othello}~所示。研究者训练了一个小型 Transformer，唯一任务是根据黑白棋落子序列预测合法落子。模型从未被告知规则，也从未见过棋盘。尽管如此，探测器仍可从其激活中恢复完整棋盘状态；当作者编辑这一内部棋盘时，模型预测的，是在编辑后的棋盘而非实际棋盘上合法的落子。然而，这项研究采用的是非线性探测器，因此第~\ref{sec:levels}~节的第一项顾虑仍然适用：也许棋盘是探测器计算出来的，而不是模型计算出来的。\citet{nanda2023emergent} 解决了这一问题。结果表明，棋盘状态可以线性解码，也可以沿线性方向编辑，并对落子产生同样的影响；但前提是将格子编码为{\em 我方（mine）}或{\em 对方（yours）}，而非黑色或白色。在这一限定下，该案例在三个层面都提供了证据。更普遍的启示还在于，模型采用的变量，未必就是我们预期的那些变量。

\begin{figure}
\centering
\newcommand{\cogminiboard}[3]{
\begin{scope}[shift={#1}]
\fill[black!8] (0,0) rectangle (1.2,1.2);
\draw[step=0.3, black!55] (0,0) grid (1.2,1.2);
\foreach \x/\y/\c in {#2} \filldraw[fill=\c, draw=black] (\x*0.3+0.15,\y*0.3+0.15) circle (0.1);
#3
\end{scope}}
\begin{tikzpicture}[font=\small,
box/.style={draw, rounded corners=2pt, align=center, text width=2.15cm, minimum height=1.05cm, inner sep=3pt},
lvl/.style={font=\footnotesize\itshape, text=black!65, align=center}]
\node[box] (seq) at (0,0) {落子序列 \texttt{\scriptsize F5 D6 C3 \dots}};
\node[box] (low) at (3.0,0) {Transformer 层 $1\dots\ell$};
\node[box] (high) at (6.4,0) {Transformer 层 $\ell{+}1\dots L$};
\node[box] (out) at (9.4,0) {预测的\\合法落子};
\draw[-latex] (seq) -- (low);
\draw[-latex] (low) -- node[above, font=\footnotesize] {$h_\ell$} (high);
\draw[-latex] (high) -- (out);
\node[lvl] at (9.4,-1.0) {(1) 行为层面};
\cogminiboard{(2.6,-4.2)}{1/1/black, 2/2/black, 1/2/white, 2/1/white, 2/3/black}{}
\cogminiboard{(5.6,-4.2)}{1/1/black, 2/2/black, 1/2/white, 2/1/white, 2/3/white}{\draw[red!70!black, very thick] (0.6,0.9) rectangle (0.9,1.2);}
\node[font=\footnotesize] at (3.2,-4.5) {解码的棋盘};
\node[font=\footnotesize] at (6.2,-4.5) {编辑后的棋盘};
\draw[-latex] (4.5,0) -- node[left=3pt, lvl, pos=0.55] {(2) 表征层面：探测} (3.2,-2.9);
\draw[-latex] (3.95,-3.6) -- node[above, font=\footnotesize] {编辑} (5.45,-3.6);
\draw[-latex] (6.2,-2.9) -- node[right=3pt, lvl, pos=0.45] {(3) 机制层面：修补编辑后的 $h_\ell$} (4.9,0);
\end{tikzpicture}
\caption{Othello-GPT 实验示意图 \citep{li2023emergent, nanda2023emergent}，该实验在第~\ref{sec:levels}~节的三个层面均提供了证据。(1) 模型仅用落子序列训练，预测合法落子。(2) 探测器从中间层的激活 $h_\ell$ 中解码棋盘状态。(3) 编辑解码出的棋盘（此处翻转一枚棋子），并相应改变激活后，模型会预测在编辑后棋盘上合法的落子。为便于阅读，图中画成 $4\times4$ 棋盘；实际棋盘为 $8\times8$。}\label{fig:othello}
\end{figure}

在游戏棋盘这一特定情境之外，\citet{gurnee2024language} 在 Llama-2 中发现了空间与时间的线性表征：城市和地标的坐标，以及历史人物与事件的日期，都可以从激活中解码出来，而且在不同提示下都具有稳健性。

在提前规划方面，对实际部署的 Claude 3.5 Haiku 模型进行的归因图分析（第~\ref{chap:mi}~章），为 Craik 所说的\emph{行动前尝试不同选择}的有限形式提供了一些证据 \citep{lindsey2025biology}。模型在写一组押韵的两行诗时，还没有写第二行，就会在第一行结尾处激活候选韵脚词的特征。抑制预定词，或注入另一个词的干预，都会改变第二行。因此，模型会预先表征一个目标状态，并利用它塑造中间生成的内容。虽然这种规划只跨越一行诗的长度，但它仍然提供了反对“下一词元预测器只能向前看一个词元”这一观点的证据。这也可能与空间规划相关，但模型是否确实会这样做，以及能做到何种程度，仍需未来研究确定。

\subsection{反对的证据}
其他研究指出了大语言模型表征能力的局限。CogEval 是一套受认知科学启发、系统评估大语言模型认知能力的规程，它测试模型处理路线、捷径与绕路的能力 \citep{momennejad2023cogmap}。当时包括 GPT-4 在内的前沿模型都出现了系统性失败：它们会虚构路线选项，并陷入循环。在使用最小对立项测试基本物理与社会世界知识的 EWoK 上，大语言模型的表现也低于人类 \citep{ivanova2025elementsworldknowledgeewok}。Citet{yamada2024spatial} 的研究表明，正确率随所需导航的结构类型而显著变化。当网格、六边形与三角形网格、环或树被描述为文本，并要求模型在其中导航时，大语言模型似乎能够较好地把握某些空间属性，却在另一些属性上失败。加入视觉似乎也无济于事：视觉语言模型（vision-language model，VLM）连 {\em left of} 与 {\em under} 这样基本的关系都难以处理 \citep{kamath2023whats}，这与视觉通道对心智理论的帮助有限这一发现相符（第~\ref{sec:theoryofmind}~节）。近期一项大规模研究表明，已记录的许多视觉语言模型问题至今依然存在 \citep{wang2026monos}。

用曼哈顿出租车行驶数据训练的生成模型，几乎可以完美预测任意路线上的下一次转弯。然而，一旦加入完全合法但不太可能出现的绕路，模型表现就会崩溃。研究者根据模型预测重建街道地图时，发现其中充斥着不可能存在的街道与立交桥 \citep{vafa2024eval}。这表明，下一词元预测的正确率，可能会让人高估模型的世界模型究竟有多连贯。在后续研究中，\citet{vafa2025foundation} 表明，用行星轨道数据训练的模型能够准确预测轨迹，但从中提取出的隐含力学定律，并不能像 Newton 定律那样很好地泛化到新的物理任务。类似地，\citet{nikankin2025arith} 表明，大语言模型依靠{\em 一组启发式规则（bag of heuristics）}而非算法与定律进行算术运算。他们发现了许多局部表征存在的证据，每个表征各自处理一种范围狭窄的数值模式。



\subsection{智能体视角}
系统层面的视角指明了一条前进道路。智能体无须将整个世界模型都承载在激活中：它可以维护一张显式地图、一个笔记文件，或一个模拟环境的程序，并在需要时查阅。当必须长期追踪空间状态时，智能体仍然举步维艰；案例研究 2 中，你已经在 BALROG 的 NetHack 任务里体验到了这一点。

此时，世界模型成为系统而非模型的属性，而且可以被检查。一个互补的发展方向，是使用从视频中学习的世界模型，生成可供智能体训练的交互式环境 \citep{bruce2024genie}。这与第~\ref{sec:openworld}~节的开放世界环境相联系。

\subsection{结语}
Othello-GPT 与曼哈顿实验的结果，如何能够同时成立？一个重要关键是训练数据的覆盖范围。在 \citet{vafa2024eval} 的研究中，用\emph{随机}生成的对局训练的黑白棋模型通过了连贯性测试，而用人类锦标赛对局训练的模型却没有通过。随机对局很可能覆盖状态空间的各个角落，而人类对局则像合理的出租车路线一样，只覆盖其中的一小部分。数据在多大程度上迫使模型恢复世界的结构，模型就能恢复到多大程度。

区分弱意义与强意义上的世界模型，有助于理解这一问题。在弱意义上，只要世界状态能够从模型激活中解码出来，模型就具有世界模型。在强意义上，这个模型还必须具有\emph{连贯性}（一个一致的模型，而非许多局部碎片），\emph{在因果意义上得到使用}，并支持向新情境泛化。

现有证据主要表明，大语言模型具有弱意义上的世界模型。取决于某个领域中的数据密度，它们也可能恢复出连贯的世界模型，达到强意义的标准：用足够多随机黑白棋对局训练的模型就是一例。另一些模型使用反映人类如何在世界中行动、如何书写世界的数据训练，而这些数据仅覆盖人们认为值得描述的部分；它们最终形成的更像是一组局部相互矛盾的草图，而非完整地图。从物理模拟中提取的规律也是如此：它们足以进行局部预测，却不能很好地泛化到其他情境。

\section{社会认知与心智理论}\label{sec:tomcognition}\label{sec:tomcog}
正如第~\ref{sec:theoryofmind}~节所见，前沿大语言模型在许多心智理论（ToM）测试上已达到人类水平。然而，测试形式稍作扰动，这种表现就会变得脆弱；用语言预测他人心理状态（字面心智理论），也并不总能转化为与预测一致的行为（功能心智理论）。这是否意味着可以说大语言模型\emph{具有}心智理论？这种心智理论是否在\emph{任何人类意义上}成立？此前我们暂时搁置了这些问题，现在将在此继续讨论。

\subsection{表征与机制层面的信念}
\citet{zhu2024language} 向两个相对较小的模型（Mistral-7B 与 DeepSeek-7B）提供错误信念故事，并基于注意力头的激活训练线性探测器。探测器能够分别解码主人公的信念，以及全知读者所知的世界状态，说明模型将两种视角区分开来。应用于 Maxi 的故事，这意味着 Maxi 相信巧克力仍在蓝色柜子里，而全知读者的信念则在巧克力趁 Maxi 不在时被移动后，更新为绿色柜子。因此，我们在第~\ref{sec:empathy}~节讨论共情时提到的自他区分，似乎在受测模型中有所表征：它们能够领会，他人（主人公）的信念，可以不同于自身（作为测试的全知读者）的信念。

在他们的实验中，沿探测方向引导激活，会改变模型在心智理论任务上的表现。在不同于训练探测器所用任务的其他任务中，也出现了这一现象，说明所表征的信念在行为中具有因果重要性。换言之，这里不仅有表征层面的证据，还有机制层面的证据。

\citet{bortoletto2024benchmarking} 将这一研究方向扩展到多个模型家族和规模。他们进一步发现，信念表征的质量与可靠性会随模型规模和指令微调而提高。他们也表明，行为层面观察到的脆弱性，同样延伸到了表征层面：全知视角的表征相当稳定，但他人信念的表征对提示变化十分敏感，甚至会受本不应产生影响的变化干扰。

\subsection{人类与机器通向相似行为的不同路径}
儿童在 3 至 7 岁间发展出心智理论。这一过程部分依赖与他人共情的先天倾向，部分依赖从环境中学习。共同注意、假装游戏、针对他人愿望和想法的对话，以及其他具身和交互经验，也都有助于心智理论的发展 \citep{apperly2011mindreaders}。语言输入同样在这一发展过程中起着重要作用。研究已经表明，心智理论的发展与认知动词的习得 \citep{milligan2007language}、否定的掌握 \citep{devilliers2014role}，以及句子补足结构的掌握 \citep{lohmann2003role} 存在联系。长期以来的大量研究尤其表明，叙事能力，即讲述和理解故事的能力，为儿童、青少年与成人心智理论技能的发展提供支撑，\citep[例如]{hutto2008folk, vanduijn2016, nicolopoulou2017narrativity}。大语言模型解决心智理论任务的能力，几乎完全依靠语言输入获得。\footnote{后训练在心智理论发展中的作用，目前大体上仍不清楚，但已有迹象表明，监督微调可以改善心智理论测试表现 \citep{kouwenhoven2026traces}。}它们用规模弥补了缺乏真实社会环境中交互嵌入的不足：大语言模型需要多出几个数量级的语言输入，心智理论能力才开始涌现 \citep{trott2023do, kouwenhoven2026traces}。这表明，即使儿童与大语言模型取得相似的测试得分，也不能据此断言它们背后的计算过程相似 \citep{hu2025reevaluatingtheorymindevaluation}。

还应补充一点：关于\emph{人类}如何处理心智理论任务的争论，本身也未有定论。按照 (1)~\emph{理论论（theory-theory）}，我们将日常心理学理论应用于他人 \citep{gopnik1992child}。按照 (2)~\emph{模拟论（simulation-theory）}，我们以自己的心智与身体为模型：想象自己置身他人的处境，以推理出他们的感受和想法 \citep{goldman2006simulating}。而 (3)~双系统观点的支持者则认为，两者缺一不可：模拟用于快速、自动的推理，理论论则补充较慢、更为审慎的社会推理 \citep{apperly2009humans}。

这三种观点是否有助于描述大语言模型中的心智理论处理，目前仍不清楚。模拟预设了一个身体，以及主体自身的心理和情绪生活，因此很难想象它在大语言模型中会是什么样子。另一方面，从第~\ref{sec:consciousness}~节将讨论的角色扮演解释来看，大语言模型似乎经常\emph{以某个角色的身份}生成文本。这与模拟存在相似之处：人类采用某种视角，生成处于那个位置的人可能会说的话。至于日常心理学理论，则可以建立在大语言模型训练数据中数量充足的文本来源之上。人类文本包含大量描述心理状态的语言。以儿童故事为例，《小红帽》（Little Red Riding Hood）的核心，围绕一个错误信念情节展开：她\emph{不知道}狼已经吃掉外婆，并躺在床上冒充外婆。小红帽\emph{纳闷}，外婆的耳朵为什么这么大，等等。从某种意义上说，这类童话，以及小说、电影剧本和其他虚构体裁，都可以视为大语言模型的大型心智理论训练集。从这些体裁中提炼出的抽象模式，很可能在大语言模型完成心智理论任务的能力中起着重要作用，但具体机制尚待揭示。

总之，对于人类，究竟是理论驱动还是模拟驱动心智理论，这场争论尚无定论；而这一区分如何适用于大语言模型，就更加不清楚。

\subsection{重访字面—功能鸿沟}
第五部针对字面心智理论与功能心智理论之间，或知与行之间的鸿沟，提出了三种候选解释：(1) 缺乏相关能力；(2) 能力存在，但不会自发运用；(3) 表征存在，但未与动作选择相连接。社会思维链（Social Chain-of-Thought）等支撑方法取得的成功，似乎支持解释 (2) \citep{akata2025playing}。这与上述探测结果相符：相关信念表征确实存在于模型中；但对于 (3)，即这些表征是否实际用于行动，目前还不能据此作出判断。

一种可能的解释是，大语言模型只有在某个情境\emph{看起来像心智理论问题}时，才会调用相关知识。如果这种要求保持隐含，相关表征就不会用于动作选择，除非明确提出问题，如社会思维链所做的那样。一些实验支持了这一解释：改变心智理论情境可能触发错误结果。例如，某个情境看起来像经典的错误信念场景，实际却只是一个简单的\emph{正确信念}案例，模型仍可能倾向于套用错误信念推理 \citep{ullman2023large, kouwenhoven2026traces}。需要注意，人类也难以跨越知与行的区别。一个人可能知道国际象棋的规则，但真正运用这些规则，在脑海中推演几步棋，需要投入大量心力，只有少数人能够掌握。Kahneman 所说的系统 2 思维，对人类而言需要高度集中注意力 \citep{kahneman2011thinking}。外部工具可以提供帮助：人类可以用一张纸追踪棋盘状态；类似地，大语言模型也可以通过控制框架增强，在外部追踪状态，正如第二部和第三部所示。

对于智能体大语言模型，决定性的实验应当在 Hanabi 一类游戏的交互\emph{过程中}，在机制层面对模型所表征的信念施加因果干预。在这种情境下，模型需要进行心智理论推理，却不会被明确问及心智理论问题。据我们所知，这样的实验尚未开展，表~\ref{tab:levels}~中心智理论对应的智能体单元格为空，也反映了这一点。

\subsection{读懂大语言模型的心思}
反过来，人类也会将自己的心智理论用于大语言模型。一项采用美国代表性样本的研究表明，多数参与者会赋予 ChatGPT 某种程度的意识体验与日常心理学意义上的体验 \citep{colombatto2024folk}。这种归因与使用频率呈正相关。这符合前文讨论的 CASA 范式（第~\ref{sec:politeness}~节），我们将在第~\ref{sec:consciousness}~节再次讨论这一点。

\subsection{结语}
根据现有证据，大语言模型的心智理论在静态测试上，可以达到与人类相当的行为与表现。有限的证据表明，这种相似性也延伸到了表征与机制层面。大语言模型将信念与现实分开表征，这些表征可以在回答心智理论问题时发挥因果作用。与人类相比，信念表征的获取和应用方式，似乎既有相似之处，也有差异。表征与行动规划之间的连接十分脆弱，往往依赖显式激活。进一步限制人类与大语言模型之间类比讨论的，是关于人类如何习得和处理心智理论的争论，本身也尚未解决。

\section{元认知、内省与自我知识} \label{sec:metacognition}
到目前为止，我们询问的是大语言模型如何表征语言、世界与他人的心智。现在要问的是，它们如何表征自身的内部运作。这对于智能体大语言模型尤其重要。为了执行多步任务，智能体必须知道，例如，自己不知道什么、何时该验证某条信息、一项任务可能需要投入多少精力，以及任务何时才真正完成。这类问题传统上都归入\emph{元认知（metacognition）}的范畴。

在心理学中，元认知通常被理解为\emph{监测}与\emph{控制}之间的循环 \citep{nelson1990meta}：系统评估自身认知过程的某个方面，例如不确定性或子任务进度，再用这一评估调节下一步行为，比如转换策略、检查中间结果、寻求帮助或停止。\citet{liu2026metacog} 在大语言模型元认知的综述中，将这一经典循环用作概念框架。这一区分与本节有关，因为支持循环两部分的证据并不一样强：模型可能表征了关于自身答案质量的信息，却不能可靠地据此行动——这又是知行鸿沟的一个例子（第~\ref{sec:tomreasoning}~节）。

还需说明术语的用法：在大语言模型文献中，\emph{元认知}一词被用于广泛的情境，从经过校准的置信度分数，到对给定输出用语言表达的反思，与心理学概念的对应都较为宽泛。因此，对于下文每项发现，我们都会询问：究竟监测了什么，我们如何知道这一点，以及监测结果是否被用于控制后续的生成步骤或动作。

\subsection{监测：知道自己知道什么}
应当区分置信度判断的两种属性。元认知\emph{校准（calibration）}关注置信度是否与正确率一致：以 80\% 的置信度给出的答案，是否也有 80\% 是正确的？元认知\emph{敏感性（sensitivity）}关注置信度能否区分正确与错误答案：模型答对时，是否比答错时更加自信？心理学家使用源于信号检测理论的测度，例如 meta-$d'$ 及比值 meta-$d'/d'$（M 比率），来表达相对于任务表现的元认知敏感性 \citep{maniscalco2012signal}。如今，这些测度也开始用于大语言模型。

早在 2022 年，\citet{kadavath2022language} 就表明，大语言模型能够相当准确地评估自己知道和不知道什么。面对各种多项选择题，模型能够给出校准良好的分数。模型也能较好地评估自身答案的正确性；还可以通过训练，让模型在答案选项出现之前，就预测自己是否知道某个问题的答案。

此后，多项研究复现并扩展了这些发现。不过，有三点需要谨慎对待。不同研究的结果并不一致，说明元认知敏感性依赖任务和领域 \citep{liu2026metacog}。第一，元认知校准与敏感性，会随任务、领域和模型而变化 \citep{liu2026metacog}。第二，采用人类反馈的后训练降低了 GPT-4 答案概率的校准程度 \citep{openai2024gpt4}，可能是因为这类训练使用的奖励模型往往偏好自信的回答 \citep{leng2025taming}。反过来，后训练也可能产生积极作用 \citep{lindsey2025introspection}。第三，结果取决于如何获取置信度：用语言表达的置信度估计、自我报告的数值（例如，提示模型在 0-100 的量表上给出置信度分数），以及根据词元概率导出的置信度，测量的并不是同一件事 \citep{liu2026metacog}。当模型被要求用语言表达置信度时，它们大多会过度自信 \citep{xiong2024can}。

\subsection{内省}
人类能否以特有方式访问自身内部状态，这一问题同样适用于大语言模型：它们是否拥有某种访问途径，而看到相同行为数据的外部观察者却不具备？\citet{binder2024looking} 微调了一个模型，让它预测自身在假设情境中的行为。与另一个用第一个模型的行为数据训练的模型相比，它的预测正确率更高。这说明，第一个模型能够访问单靠行为无法获得的内部信息。相反，将模型对自身语言知识作出的语言判断，例如某个句子是否合乎语法，与其实际赋予句子的概率比较时，这些判断对自身概率的追踪，并不比对其他模型概率的追踪更好 \citep{song2025languagemodelsfailintrospect}。

\citet{lindsey2025introspection} 通过\emph{概念注入（concept injection）}，对内省进行了因果检验。例如，研究者将对应于全大写书写（ALL CAPS）的激活模式，加入模型处理简单查询时的激活。当被问及是否察觉任何异常时，至少在部分试次中，模型能够合理地用语言表达所注入的概念。在 ALL CAPS 的例子中，它可以表示自己察觉到了一个“过于强烈、音量很高的概念”。

研究者也曾询问，推理轨迹是否可以视为反映内部状态的自我报告。\citet{turpin2023language} 表明，如果提示中的多项选择题示例，总是把正确答案放在第一个位置，模型就会追随这一偏差。当提示要求解释时，它们通常会为所选答案编造合理化理由，而非说明这一偏差。这与 \citet{chen2025reasoningmodelsdontsay} 的发现一致。当提示中植入关于答案的暗示时，模型会利用它们，却仅在少数推理轨迹中承认这一点。需要注意，人类也会做类似的事情。在经典实验中，参与者自信地为自己的选择给出理由，但可以证明，这些理由并不是作出那些选择的原因 \citep{nisbett1977telling}。

总体而言，支持模型能够以特有途径访问自身内部状态的研究结果并不一致。即使这种能力存在，它也范围狭窄、不可靠，并且依赖实验设置。

\subsection{智能体视角}
\citet{shapira2026agentschaos} 研究中的智能体声称任务已经完成，系统状态却显示事实并非如此。这可以看作自我监测的失败：关于自身状态的主张，从未与世界核对。对于人类，工作与其他日常场景中的许多这类核对已成习惯；但对于智能体，必须将其设计为验证步骤、测试，或每次动作后的状态检查。智能体架构正是这样做的，例如设置独立监测器，检测深度搜索智能体的检索结果与推理何时不再一致，并触发纠正动作 \citep{sun2026deepsearch}。与世界模型一样，元认知也可能出现在整个系统层面，而非单个模型内部。

\subsection{结语}
大语言模型的元认知、内省与自我知识可以被检测到，但也在多个方面受到限制。在监测方面，模型的语言表达、内部概率与激活，携带着程度不一的答案可靠性信息。关于特有内省途径的证据并不一致，充其量只支持有限的访问能力。在控制方面，证据更加稀薄：监测信号往往没有被用于调节行为。这为智能体系统留下了设计空间，使其能够实现与人类元认知控制在功能上等效的行为。

\section{意识}\label{sec:consciousness}
机器能否具有意识？这个问题由来已久，过去属于推测、思想实验与科幻的领域。然而今天，我们看到，意识研究领域的顶尖科学家和哲学家已经开始围绕大语言模型与智能体大语言模型讨论这一问题。本节将概述可以如何研究它，以及为什么它仍未有定论。

\subsection{定义}
意识与智能这一更宽泛的概念，有何区别？在经典论文 What Is It Like to Be a Bat? 中，\citet{nagel1974bat} 提出：如果成为某个系统有“某种体验”，那么该系统就具有\emph{现象意识（phenomenal consciousness）}。这将意识定位为一种需要\emph{从内部}体验的属性，也使它按定义便很难从外部视角获知。\citet{block1995confusion} 将其与\emph{通达意识（access consciousness）}区分开来，后者定义为：来自感官、记忆及其他官能的信息，能够用于推理、反思与行动控制。这两类意识都应与\emph{感受能力（sentience）}区分，感受能力指具有某种效价的体验，例如愉悦、痛苦与欲望。空间定向、心智理论和上述元认知等功能性能力，原则上可以在没有感受能力或意识的情况下存在。反过来，虽然多数动物至少在后两种功能性能力上的得分不高，但它们很可能具有感受能力。

作为人类观察者和大语言模型用户，我们经常随意使用诸如\emph{系统想要 $\dots$}或\emph{智能体决定 $\ldots$}这样的表达。必须区分：出于语用目的使用这类归因，与承诺这些心理状态在本体论上具有何种地位，是两回事 \citep{van-dijk-etal-2023-large}。\citet{dennett1987intentional} 的\emph{意向立场（intentional stance）}也作出了这一区分：只要赋予某个系统信念、欲望和其他心理状态，能够有效预测其行为，就采取这种做法。说一个恒温器\emph{发现房间太冷就会打开暖气}，能够为这一系统提供可用的模型，而不需要任何证据，证明它具有与我们以相同方式运作的内心生活。有些人认为，\emph{大语言模型理解 $\ldots$}或\emph{这个智能体想要 $\ldots$}之类表达的含义仅止于此；另一些人则认为，用在大语言模型上时，这类表达还有值得探索的深层内涵。

\subsection{功能主义与理论驱动的评估}
与第五部和前文的讨论一致，\citet{butlin2023consciousness} 认为，仅凭行为很难了解大语言模型的内部运作。因此，他们建议从主要的科学意识理论出发，导出\emph{指标属性（indicator properties）}，再据此检验当前人工智能架构与系统属性。他们的论述建立在\emph{计算功能主义（computational functionalism）}的前提上：执行适当类型的计算，是意识的必要且充分条件。关键在于系统的功能组织，即它表征哪些信息、如何转换信息、各组件如何交互，而非实现这些处理的物质载体。

在这一观点中，意识\emph{与实现载体无关（substrate-independent）}，人们常用以下思想实验来说明这一点。假设你的脑部受伤，只损伤了一个或几个神经元，而微型医疗技术的进步，使受损神经元能够被硅基对应物替代，后者与原来的生物神经元发挥完全相同的功能。你的脑功能在功能意义上完全等效。你会说自己不再有意识了吗？如果后续受伤后，更多神经元被替换，直到四分之一的大脑由硅基神经元等效物构成呢？如果是一半，或四分之三呢？沿着 \citet{searle1980minds} 的思路，生物学取向的观点仍然质疑这一看法，认为意识依赖活体神经组织的特定属性（见下文第~\ref{againstfunct}~节）。

\citet{butlin2023consciousness} 认为，应将功能主义作为工作假设，而非既定事实，因为它使意识问题可以接受经验研究：明确指出相关计算的科学理论，能够转化为人工智能架构中可检查的属性。下面列出主要指标与作者的评估，但需注意，这篇论文写于 2023 年，缺少智能体维度；我们将在下文回到这一点。

\paragraph{循环加工} 意识与循环加工（recurrent processing）而非纯前馈加工相关。算法层面的循环，在有限意义上是循环神经网络（RNN）或长短期记忆网络（LSTM）的属性；但 Transformer 的前向传播并不具有循环性。

\paragraph{全局工作空间理论} 对于人类，全局工作空间理论（global workspace theory）认为，专门化的并行模块会竞争一个容量有限的工作空间，其内容再向所有模块进行全局广播。有人提出，Transformer 的残差流（第~\ref{residualstream}~节）发挥了这种工作空间的作用；但这一论据较弱，因为残差流并不是明确的瓶颈，也不存在真正的循环或向模块回传的广播。不过，近期研究提出，残差流中一个稀疏、容量有限的子空间确实表现为瓶颈（见下文）\citep{gurnee2026verbalize}。

\paragraph{高阶理论} 按照高阶理论（higher-order theory），当一个高阶元认知机制监测一阶知觉表征，例如将其标记为可靠信息而非噪声，并向信念形成与动作选择系统提供信息时，某个状态便成为有意识的状态。作者认为，当前没有系统实现了这一组合；但他们也强调，使用现有组件，在技术上可以做到这一点。

\paragraph{注意力图式理论} 按照注意力图式理论（attention schema theory），意识源于系统对自身注意力的简化预测模型，系统借此表征并控制自身当前的注意力状态。作者指出，机器注意力与这一理论所指的注意力，只是宽泛的类比，当然也不能等同于可用来引导当前注意力的自我模型。

\paragraph{预测加工} 在预测加工（predictive processing）理论中，知觉被表述为预测误差最小化。大语言模型的训练方式中包含这方面的要素，但作者更倾向于将其视为一般性的认知框架，而非独立的意识指标。

\paragraph{能动性} 按照这一观点，从反馈中学习，并根据具体目标选择输出，被视为意识的一种形式。作者提出，强化学习足以实现其基本形式；但它涉及的是外部设定的目标，这些目标本身并不是系统的一部分。

\paragraph{具身性} 作者强调，关键并不在于物理身体本身，而在于拥有一个关于系统输出如何影响其输入的模型，这有助于区分由自身引起和由世界引起的变化。他们讨论了多种具备这类模型变体的系统，有些通过控制机器人，有些通过控制虚拟化身来实现；但作者认为，对于动物界普遍具有的、完整且端到端的具身性而言，这些系统仍然只是依据较弱或尚不确定的候选者。

总体而言，作者的立场是：按照他们定义的指标，当前没有任何人工智能系统是具有意识的有力候选者；但他们也看不到明显的技术障碍，会阻止人们在不久的将来构建满足其中许多指标的系统。正如下文将看到的，智能体维度与这场争论密切相关。

\citet{chalmers2023could} 在同一年得出了类似结论。他列出了自己认为大语言模型缺少的要素：生物学基础、感官、身体、世界模型与自我模型、循环加工、全局工作空间，以及统一的能动性。其中两项障碍正是上述第~\ref{sec:worldmodels}~节和第~\ref{sec:metacognition}~节的主题；那里的证据说明，这些障碍的多个方面目前已被克服，或正在发展之中。Chalmers 认为当前大语言模型具有意识的概率低于 10\%，但对经过扩展的后继系统给出了高得多的概率。表~\ref{tab:indicators}~提供了简化概览。

近期的机制可解释性研究，开始关注前沿大语言模型的前向传播中是否会涌现全局工作空间。\citet{gurnee2026verbalize} 提出了\emph{雅可比透镜（Jacobian Lens）}：为词表中的每个词元，在残差流中指定一个方向，用以编码模型将该词元表达出来的倾向，并在不同上下文间取平均。他们识别出\emph{J 空间（J-space）}，它同一时刻只容纳几十个活跃概念；他们认为，它具有全局工作空间理论赋予意识通达的功能特征。其内容在很大程度上决定模型会报告什么，可以在模型做其他事情时按指令调取，承载多步推理中未说出口的中间步骤，还能通过替换，一次性改变许多不同下游计算的方向。消融这一空间后，句法分析、分类与流畅续写等常规处理大体保持完好，但会损害灵活推理；正是这类推理，使当前前沿模型在一系列复杂任务中具有强大能力。作者并不对现象意识表态，并强调这里的相似是功能层面的，而非架构层面的：没有来自封装模块的输入，也没有向它们返回的反馈；广播沿单次前馈传播的层深展开，而非通过循环回路进行。这些发现加强了 \citet{butlin2023consciousness} 所倡导的基于指标的方法，却没有触及与生物自然主义者的分歧（见下文第~\ref{againstfunct}~节），因为这一分歧不在于关于大语言模型的事实，而在于对“具有意识”意味着什么的前提。

\begin{table}
\footnotesize
\begin{tabularx}{\textwidth}{@{}p{2.6cm}p{3.4cm}XX@{}}
\toprule
\textbf{理论} & \textbf{指标（简化）} & \textbf{未经额外增强的大语言模型} & \textbf{智能体系统} \\
\midrule
循环加工 & 处理中的反馈回路 & 有争议：每个词元的处理为前馈，仅通过自身输出形成循环 & 相同 \\
全局工作空间 & 容量有限、具有全局广播的工作空间 & 有争议 & 更为接近：共享记忆缓冲区向感知、规划与动作模块提供信息 \\
高阶理论 & 监测自身表征 & 较弱且不可靠（第~\ref{sec:metacognition}~节） & 验证模块监测输出，而非内部状态 \\
注意力图式 & 关于自身注意力的模型 & 没有证据 & 没有证据 \\
能动性与具身性 & 利用反馈追求目标；身体模型 & 大体缺乏 & 部分具有能动性；具身性仅见于机器人领域 \\
生物学观点 & 生命、代谢、身体调节 & 缺乏 & 缺乏 \\
\bottomrule
\end{tabularx}
\caption{意识指标属性及未经额外增强的大语言模型与智能体系统满足这些指标程度的简化概览。依据 \citet{butlin2023consciousness} 和 \citet{goldstein2024case} 整理；最后两列的评估仍有争议。}\label{tab:indicators}
\end{table}

\subsection{对功能主义的反对} \label{againstfunct}
\citep{seth2025conscious} 反对功能主义，为一种称作\emph{生物自然主义（biological naturalism）}的立场辩护。按照这一立场，意识与我们作为生命体的本性紧密交织：它的功能，是调节一个处境脆弱、必须维持自身存在的身体。预测误差最小化等智能系统属性，以及空间定向、心智理论和自我反思等能力，都应被理解为服务于生存。在自然系统与人工系统中，这些属性和能力可能存在局部相似；但在人工系统中，它们并未以因果方式连接在一条生物性的存在链条上。换言之，在生物自然主义者看来，具有意识就需要承受切身的利害后果。\footnote{Seth 的文章发表于 \textit{Brain and Behavioural Sciences}，按照该期刊的惯常形式，文章后附有公开同行评论。Leonard Dung 的 Why I Am Not a Biological Naturalist 是其中一篇评论，包含一项反驳。}\citet{aru2023feasibility} 赞同这一论证思路，并补充说，大语言模型的输入缺少我们通过感官接触世界时所具有的丰富 4E 信息内容（见第~\ref{sec:unit}~节）。因此，从神经科学视角看，大语言模型缺少丘脑皮层系统的一些关键特征，而这些特征已知与哺乳动物的意识相关。

对于生物自然主义者，随着系统 (1)~变得更加类似生命，并开始拥有自主的演化与发育轨迹，且其智能属性与这些轨迹存在因果联系；(2)~拥有与外部世界连接的更丰富信息流，人工意识就更有可能出现。因此，应当注意，他们与计算功能主义者的分歧，主要在于意识是什么这一前提，而非关于大语言模型的事实。

\subsection{自我报告与内心生活}\label{sec:selfreports}
大语言模型会以第一人称报告体验、思考步骤、情绪等内容。至少有两个理由表明，不应对这些报告赋予太大权重。第一，大语言模型用人类文本训练，而这些文本中充满了这类描述。因此，这些描述未必与大语言模型自身及当前上下文相联系，也可能独立于两者之外。第二，在 \citet{shanahan2023roleplaylargelanguagemodels} 的\emph{角色扮演（role-play）}解释中，模型在结构上就是作为某个角色行动的。这同样使自我报告与模型本身脱钩：同一个模型既可以扮演声称自己具有意识的角色，也同样可以扮演否认自己具有意识的角色。这就是所谓的\emph{操纵问题（gaming problem）}：训练数据中表征心理属性的标志过于丰富，致使这些标志无法再被可靠使用 \citep{birch2024gaming}。有人提出，如果训练模型针对那些答案可以通过计算验证的问题进行准确内省，自我报告可能会变得更有信息量、更可靠 \citep{perez2023eval}；但第~\ref{sec:metacognition}~节表明，目前我们距离这种情况还有很远。

在第~\ref{sec:empathy}~节，我们搁置了大语言模型是否有所\emph{感受}的问题。在行为层面，诱发焦虑的提示会改变后续自我报告：它们使大语言模型在焦虑问卷上的得分升高，而放松提示则会部分逆转这一效应 \citep{BenZion2026anxiety}。上述两点顾虑同样适用于这些发现：模型可以依靠训练中学到的、焦虑不同方面之间的关联来回答问题，而无须涉及其实际内部状态。在表征层面，情绪概念对应于激活空间中的方向；在机制层面，沿这些方向引导激活会改变模型输出 \citep{zou2023representation}。因此，可以认为，这些状态调节行为的方式，与情绪在人类和其他动物中的作用颇为相似。按照功能主义的解释，这使它们成为类似情绪的状态 \citep{Moerland2018emotion}；但这些状态是否伴随着某种可与人类和其他动物的情绪体验相比的感受，仍然存在上述全部不确定性。

\subsection{智能体视角}
在撰写本书时，专门针对意识与智能体大语言模型交叉领域的研究仍然稀少。不过，已有论证指出，智能体架构至少与一种大语言模型意识理论相关。\citet{goldstein2024doeschatgptmind} 观察到，第~\ref{chap:emergent}~章讨论的那类大语言模型智能体，具有记忆流，以及通过共享缓冲区通信的检索、反思和规划模块，已接近全局工作空间理论（Global Workspace Theory，GWT）要求的架构。他们认为，如果 GWT 正确，那么这些智能体已经具有意识，或很容易被改造成具有意识的系统；如果有人认为这一结论不可信，就应将其视为反对 GWT 的证据。持久记忆与长期目标，也回应了 Chalmers 指出的缺乏统一能动性的问题。

\citet{gurnee2026verbalize} 提出了一点有趣的观察：J 空间可能在一次前向传播中，模拟某些也被归于智能体架构和思维链提示中循环加工的功能。通过写出推理过程而解决的问题，对 J 空间消融的抵抗力要强得多；这说明，在没有写出过程时，模型本来会在内部承载这类推理。这使争论重新回到一个问题：任何形式的认知功能，究竟应归于哪个层面？严格发生于模型内部的过程，与系统层面涌现的过程，其边界无法先验划定，而取决于任务、提示格式、工具是否可用、交互以及其他情境因素。需要注意，认知外在主义者对人类认知功能提出的，恰恰也是这一观点 \citep{noe2009outofourheads}。

\subsection{结语}
对于意识与内部状态体验，本章开头介绍的两种潜在错误，都可能代价高昂。拟人化已经在发生，并带来了情感依赖、信任错置与操纵的风险。如果这类系统在某种程度上具有意识或感受能力，类人性否认则会带来大规模的道德错误。\citet{birch2024gaming} 主张采取与相关概率相称的预防措施；\citet{butlin2023consciousness} 则认为，概率已经高到足以立即着手评估系统。有些开发者也确实启动了模型福祉研究计划。例如，Anthropic 在 Opus 4 模型发布文档中加入了福祉评估，并允许部分模型结束持续受到辱骂的对话，将其作为一种低成本的预防措施，以应对模型福祉可能确实重要的情形。\footnote{\url{https://www.anthropic.com/research/end-subset-conversations}}无论如何看待这些概率，这一主题都已经从哲学研讨课走进了工程实践。


\section{结语}
本章询问：第五部所综述的行为背后是什么，它是否类似于驱动人类与动物行为的认知？我们由此梳理出四条总体线索。

第一，内部表征以人类观察者能够理解的方式组织到何种程度，会随任务和领域而变化。用文本训练的大语言模型，可以形成映射地理、棋盘状态、故事人物信念，以及自身所扮演角色特征的表征；在某些情况下，还可以通过机制研究证明，这些表征确实被使用。随机鹦鹉的形象并不符合这些发现：毕竟，只要至少在某些情况下存在从表征到行为的因果链条，就足以否证“大语言模型仅在表面层次运作”的主张。然而，同样的方法也揭示，表征可能多么不连贯，回应又可能多么依赖零散的启发式。这部分源于人类文本对世界的覆盖并不均匀，部分源于在某些领域，表面统计特征就足以满足训练目标。因此，一个能够\emph{理解}的心智形象也并不适合。大语言模型是一类新的系统；我们为人类和其他动物发展出的认知词汇，应用于它们时只能部分贴合。

第二，知行鸿沟再次出现在认知领域。模型表征的信息，多于它们自发使用或报告的信息：我们看到了不指导行动的信念状态、未被应用的定义，以及被错误描述的计算。一种合理假设是，结构化提示与支撑架构之所以有效，是因为它们将现有表征与行为连接起来。智能体循环在理论上能够自主提供这种结构和支撑。目前，这一点几乎还没有接受经验检验：所有表征和机制证据都涉及未经额外增强的模型。对于智能体，表~\ref{tab:levels}~的最后一列只有行为条目。


第三，现在我们可以将智能体扩展描述为\emph{将认知从模型移向系统}。世界状态可以保存在显式、外部存储且持续更新的地图中；自我监测可以通过验证步骤实现；智能体之间共享的记忆，则可以充当工作空间。按照哲学中的\emph{延展心智（extended mind）}论题，如果我们可靠地查阅一个笔记本，将其作为外部记忆，那么对于人类，它也会成为认知系统的一部分 \citep{clark1998extended}。智能体大语言模型的认知架构，正是遵循这一思路设计的 \citep{sumers2024cognarch}。其一个优点是，外部化的认知可以被检查，因此提高了这项技术的透明度。

第四，对于意识，诚实的立场是在各个层面都承认不确定性。不同理论对于任何数字系统能否具有意识存在分歧；自我报告与其他行为特征并不可靠，无法据此判定；目前也不清楚，如何让可解释性技术为这一问题提供有意义的信息。现在可以做的，是避免那些会误导用户、使其过度赋予系统意识的设计。研究者认为，现在应当开始评估模型福祉 \citep{butlin2023consciousness,long2024taking}。


\bibliographystyle{plainnat}
\bibliography{agenticbook,refs,case,synthetic}
